\section{Experiments}
\label{sec:discovery_agents}
We benchmark state-of-the-art LLM-based few-shot reasoning methods as discovery agents with two closed models, GPT-4o and GPT-4-0125-preview (GPT-4p), and one open, Llama-3-70B, model powering the reasoning methods. A discovery agent takes the task description, paths to the task dataset(s) $D$, metadata about the datasets (description, column descriptions), and the goal, $G$, to produce a natural language (NL) hypothesis specified by context, variables, and relationship. 

% \textbf{Output format:} an NL hypothesis + (optional) used workflow in NL

% \begin{itemize}[left=2pt, itemsep=0.1pt]
\begin{ite}
    \item \textbf{CodeGen} generates the entire code at one go to solve the task, where we provide a demonstration of a solution code in the context. After code execution and based on the result, it generates the NL hypothesis and summarizes the workflow.
    \item \textbf{ReAct} \cite{yao2023react} solves the task by generating thought and subsequent codes in a multi-turn fashion.
    % code-generation and ability to solve the task. For each turn, the agent generates a "think" statement followed by a Python code. Upon executing code, the system performs a think $\rightarrow$ code generation loop till it thinks it has solved the task. Otherwise, we cut off after a fixed step limit.
    \item \textbf{Reflexion} \cite{Shinn2023ReflexionLA} is an extension of CodeGen agent, where at the end of one trial, we provide the oracle HMS score as an evaluation signal, and it generates a reflection to improve (when $\mathrm{HMS} < 1$) in the next trial till it solves the task, or maximum trials (3) are reached.
    % We extend the Base agent with reflection capability. After one (only) step of code execution and final hypothesis generation, we allow the agent to reflect if the answer is not correct (oracle reward). Then, we restart the agent to solve the task again, but along with the input, the additional reflection is provided to the agent.
    \item \textbf{DataVoyager} is a multi-component data-driven discovery agent, introduced in \cite{majumder2024data}. It has four components, planner, code generator, data analysis, and critic, that orchestrate the discovery process.
    % \item \textbf{ReAct + Reflexion:} Now we replace the Base agent with the ReAct agent in the previously discussed Reflexion agent---this now has the ability to think while solving the task and reflect at the end of the task.
\end{ite}
% \end{itemize}

% \textbf{Experiments}
% Results from the running discovery agents on DiscoveryBench.

% \textbf{API-based:} GPT-4o, GPT4-0125-preview, \textbf{Open} Llama-3


% \textbf{Zero-shot vs. Few-shot:}

% \subsection{Single-agent Frameworks}
% Supported models: Claude, Llama, Gemma, OLMo, Gemini (with diff features like code-llama); Supp. prompts: ReAct, Plan-and-Execute, Retry


% \subsection{Multi-agent Frameworks}
% Supported models: GPT-4, 3.5, turbo; Supp. prompts: ReAct, Plan-and-Execute, Retry